% Source: https://artofproblemsolving.com/community/c3h3166494_2021_lmt_fall_guts_round_p13p24_lexington_mathematical_tournament
% Stable identity: aops-topic-3166494-q16
The first move must increase the radius to one. At any fixed positive radius, the walker can make $0$ through $359$ angular moves before moving outward or stopping; a 360th angular move would revisit its arrival point on that circle. Radii never decrease, so these are the only revisit restrictions. A path ending at radius $r$ is uniquely determined by one of $360$ angular-run lengths at each of the $r$ radii. Summing over all endpoint angles therefore gives $360^r$ paths. The requested total is $\sum_{r=1}^{2020}360^r=\boxed{\dfrac{360(360^{2020}-1)}{359}}$.
\begin{asy}
usepackage("amsmath");
size(7cm);
for(int r=1;r<=3;++r)draw(circle((0,0),r),gray);draw((0,0)--(1,0),blue,Arrow);draw(arc((0,0),1,0,50),blue,Arrow);draw(dir(50)--2*dir(50),blue,Arrow);draw(arc((0,0),2,50,120),blue,Arrow);label("radial and angular moves",(0,-3.5));
\end{asy}
